\hypertarget{appendix-a59658325cd6}{}\label{referencestart:appendix-a59658325cd6}
\section{Function contracts: domains, codomains, and images}\label{appendix-a59658325cd6-function-contracts-domains-codomains-and-images}

\textbf{Curriculum depth:} Developed · \textbf{First seen:} Phase 9 · \textbf{Latest development:} Phase 9

\subsection{Definition and mechanism}\label{appendix-a59658325cd6-definition-and-mechanism}

A function is a rule that assigns exactly one output to every allowed input. Its domain states which inputs are valid; its codomain states the declared output set; its image is the subset actually reached. These sets affect whether an inverse exists and whether a formula is defined at a boundary. For example, squaring maps the real line into nonnegative outputs but is not one-to-one until its domain is restricted to nonnegative values. In a camera map, division by depth makes zero depth invalid, while a positive-forward convention further restricts visible points. Treating these as contracts makes composition safer: the output of one function must lie in the next function's domain. A typed implementation can encode some, but not all, of these mathematical restrictions.

\subsection{Prerequisites}\label{appendix-a59658325cd6-prerequisites}

\begin{itemize}
\tightlist
\item
  No earlier concept node is required in this planned path.
\end{itemize}

\subsection{Used by}\label{appendix-a59658325cd6-used-by}

\begin{itemize}
\tightlist
\item
  \hyperref[appendix-d9f70bdfbc92]{Differentiation rules and chain rule}
\item
  \hyperref[appendix-677b59746ae0]{Partial derivatives and gradients}
\end{itemize}

\subsection{Full teaching explanation}\label{appendix-a59658325cd6-full-teaching-explanation}

\begin{itemize}
\tightlist
\item
  \href{../02-Mathematics-for-Visual-Computing/main.pdf\#nameddest=01-foundations-i-ch020-functions-limits-continuity-and-scalar-change}{Phase 9 teaching section (Mathematics for Visual Computing, Mathematics for Visual Computing, Chapter 4)}
\item
  \href{../01-Bits-Programs-and-Reliable-Systems/main.pdf\#nameddest=appendix-adf48f5dc5c1}{Visual Recall/09 - Calculus and Optimization - Visual Recall.md (Bits, Programs, and Reliable Systems, Reference notes: 09 - Calculus and Optimization - Visual Recall)} · visual recall
\end{itemize}

\subsection{Source coverage}\label{appendix-a59658325cd6-source-coverage}

Supplied lesson parts: 801 (T04 · M049).
